\authoredsection{Agile Process Implementation}{Mohsin}{Mohsin Saeed}\label{sec:agile-process}

Scrum formed the project's organizational foundation from the outset, but our process did not remain static. Across five sprints the team progressively adapted its practices to better support a five-member, multi-module development effort, and by the final sprint these incremental improvements had produced a lightweight but highly coordinated workflow. This section describes the Scrum structure we adopted, the collaboration tools that supported it, and how the process evolved across the project.

\subsection{Scrum Roles, Events, and Artifacts}\label{sec:scrum-roles}

Our Scrum setup covered the full framework, roles, events, and artifacts, which we then adapted to fit a five-person student team. Interhyp's representatives acted as Product Owners, setting direction and priorities from the domain side, while the five of us formed the Development Team. The Scrum Master role rotated through all five members, one per sprint, and in our team it carried more than facilitation. In practice, the acting Scrum Master assumed operational ownership of the sprint, managing the Jira backlog, setting up and closing each iteration, leading the review with Interhyp, and carrying the product owners' feedback into the next sprint's planning.

The events ran on a fixed two-week rhythm across five sprints, the first beginning on 28 April 2026. Each sprint opened with planning, was paced by regular stand-ups, and closed with a review alongside the product owners followed by a team retrospective. The artifacts lived in Jira throughout: a continuously refined product backlog fed a committed sprint backlog organized around the system's major components, with epics such as the router, the judge, and the optimizer broken down into user stories and story-pointed tasks, and every sprint produced a working increment, from the first runnable chatbot to the fully integrated pipeline. Burndown charts gave the team an objective, shared view of whether each sprint's commitments were on track.

Within this structure, the most visible adaptation was our meeting cadence. With five members of uneven Scrum experience, each balancing the project against jobs and coursework, a by-the-book daily stand-up would have cost more than it returned. We settled in Sprint 1 on three video stand-ups a week over Google Meet, on Tuesday, Thursday, and Sunday, with extra calls arranged as needed when work between sessions required closer coordination. Meeting less often kept overhead low and left long, uninterrupted blocks for the implementation and research the project depended on, while the regular checkpoints still kept everyone's progress visible. This cadence, however, was never treated as fixed. Instead, it became a process parameter that the team adjusted throughout the project as coordination demands evolved.

\subsection{Collaboration and Development Tools}\label{sec:collaboration-tools}

Our toolset was deliberately small, with each tool chosen to solve a specific problem for a team that was rarely in one place. Planning and work tracking centered on Jira, a choice suggested by Interhyp and aligned with common industry practice for agile software projects, which held the epics, user stories, and story-point estimates that kept scope and status transparent to everyone \autocite{avramovska2024jira}. Source control was managed through GitHub, organized under a shared LMUxInterhyp space the product owners could access and split into two repositories, one for the chatbot application and one for pipeline experiments, so that fast experimentation on the optimization modules would not slow the main application; alongside the code it held the datasets and other artifacts the pipeline needed to run. Work followed a disciplined branch-and-review workflow, with each ticket developed on its own feature branch named after the corresponding Jira ID and merged only after a pull-request review, which kept planned and delivered work traceable to one another. A shared Google Drive was the home for everything text-based the team produced and wanted to share, including reports, written documentation, and working drafts. A WhatsApp group handled quick informal coordination and confirmations between meetings, keeping lightweight communication off the formal channels. For visual and collaborative work, we used two surfaces: Miro as a digital whiteboard for remote brainstorming, mapping module contracts, and running interactive retrospectives, and a physical whiteboard during the in-person sessions at the LMU library, where design decisions were worked out together. Each tool held a clear role, and together they preserved transparency and short feedback loops despite the team seldom being co-located.

\subsection{Iterative Refinement of the Process}\label{sec:process-refinement}

The ritual that shaped how the team worked was the sprint retrospective, which the acting Scrum Master facilitated at the close of each sprint. What made these sessions matter was that they consistently turned reflection into concrete change for the next sprint, and over successive sprints they also strengthened the team's coordination and alignment, an effect \textcite{spiegler2018leadership} associate with maturing retrospective practice. The clearest example is how the team's own way of meeting evolved.

The first such change was the weekly Focus Tuesday session at the LMU library, which came directly out of the Sprint 1 retrospective. This two-hour in-person block became the place where the team settled design questions on a shared whiteboard, integrated components side by side, and brought less experienced members up to speed through pair work. Such in-person working sessions are where agile teams reach collaborative decisions most quickly \autocite{moe2012shared}, and the pair work itself is consistent with the finding of \textcite{hannay2009pair} that pairing improves solution correctness and knowledge transfer at a modest short-term cost. Decisions that would have stalled across a week of messages were routinely closed in a single session.

From there the cadence itself became something the team tuned rather than fixed. It began at three online stand-ups a week; the Sprint 3 retrospective raised it to four for Sprint 4, when the separate modules were being integrated and coordination needs peaked. During that stretch the team also arranged ad-hoc calls whenever the work demanded them, some running as late as 1am as integration problems were worked through in real time, and individual sessions were allowed to run long when a push was needed, as when a stand-up in the final sprint was extended to drive the Optimization Dashboard forward.

This is the shift that defined our process: the team started by following Scrum closely and, sprint by sprint, learned to treat its rituals as instruments to be adjusted to the work rather than obligations to be met. Meeting more when integration was intense, and stretching sessions when a deadline was near, reflects the distinction drawn by \textcite{benlian2021sprint} between sprint zeal and sprint fatigue, that agile practices energize a team only when they fit its current state, which is the fit the team kept adjusting toward.
